\title{X-LoRA: Mixture of Low-Rank Adapter Experts, a Flexible Framework for Large Language Models with Applications in Protein Mechanics and Molecular Design}

\begin{abstract}
We present a mixture of experts approach to develop fine-tuned large language models utilizing a deep, layer-wise, token-level method grounded in low-rank adaptation (LoRA). Beginning with a collection of pre-trained LoRA adapters, our gating mechanism leverages hidden states to dynamically combine adapted layers, enabling the resulting X-LoRA model to harness diverse capabilities and form novel deep layer-wise combinations to tackle tasks. This design draws inspiration from biological concepts of universality and diversity, where neural network components are repurposed across various hierarchical structures. Consequently, the X-LoRA model can be seamlessly integrated with any existing large language model (LLM) without requiring modifications to its fundamental architecture. We create a customized X-LoRA model that delivers scientific functionalities such as forward/inverse analysis tasks and improved reasoning skills, concentrating on biomaterial analysis, protein mechanics, and design. The contributions of this work include providing access to easily extendable and adaptable models equipped with substantial domain expertise and the ability to synthesize knowledge across different fields. Featuring experts in biology, mathematics, reasoning, bio-inspired materials, mechanics and materials science, chemistry, protein biophysics, mechanics, and quantum-mechanics-based molecular properties, we carry out a series of physics-oriented case studies. These studies assess knowledge retrieval, protein mechanics forward/inverse tasks, protein design, adversarial agentic modeling including ontological knowledge graph construction, and molecular design. The model is not only capable of producing quantitative predictions for nanomechanical properties of proteins or quantum mechanical molecular attributes but also performs reasoning on the outcomes and accurately predicts plausible mechanisms explaining distinct molecular behaviors.
\end{abstract}

\section{Introduction}
Large language models (LLMs)~\cite{Vaswani2017AttentionNeedc,Touvron2023LlamaModels,OpenAI2023GPT-4Report,Chowdhery2022PaLM:Pathways,Jiang2023Mistral7B,Gunasekar2023TextbooksNeed} have attracted considerable attention, including in the creation of specialized models that excel in particular tasks, reasoning processes, or scientific fields~\cite{Bubeck2023SparksGPT-4,Buehler2023MechGPTModalities,Nejjar2023LLMsAnalysis,Buehler2023GenerativeDesign,Luu2023BioinspiredLLM:Materials,Luu2023GenerativeSolvents,Buehler2023MeLMProblemsb,Ge2023OpenAGI:Experts}. Nevertheless, the training of such models can be expensive, especially when a wide range of capabilities is required. Techniques like low-rank adapters (LoRA)\cite{hu2021lora} have been introduced as a more resource-efficient alternative, but these adaptations typically focus on more specific domains of knowledge. The core idea behind LoRA is to employ low-rank matrices added to the original full-scale matrix, designating these low-rank matrices as the only components updated during training. Because only the adapter layers undergo training, these models generally retain the knowledge acquired in the pre-training phase of the base model, while making the model more suitable for particular tasks and requiring fewer computational resources. The efficiency of training LoRA adapters enables the development of a large number of specialized models using this approach.

Although LoRA adapters offer a resource-effective path to specialized capabilities, it remains unclear how to combine multiple such adapters into enhanced models that deliver a cohesive and potentially superior set of functionalities. Indeed, combining LLMs has emerged as a compelling research direction~\cite{Kim2023SOLARUp-Scaling}. Experimental efforts, such as employing spherically interpolated model parameters (SLERP) and related techniques~\cite{Arcee-ai/mergekit:Models.}, have demonstrated promising outcomes and the potential of model mixing to unlock new capabilities. These mixed models are generally generated {\it a priori} via specific algorithms and, while they have shown encouraging results on some benchmarks, further investigation is needed to systematically design mixed models and understand how certain features arise.

We propose that dynamically mixing parameters could represent a more universal method, as it enables the combined model to investigate novel combinations during inference.

Similar methodologies have employed mixture of experts (MoE) techniques~\cite{Jacobs1991AdaptiveExperts,Hampshire1992TheRecognition,Jordan1994HierarchicalAlgorithm}, including the recently introduced Mixtral LLM~\cite{Jiang2024MixtralExperts}. Nonetheless, these methods tend to be computationally intensive, requiring considerable resources even at inference time. In contrast, we present a computationally efficient mixture-of-expert strategy utilizing a collection of distinct LoRA adapters, each of which can be easily trained independently. Our method features a fine-grained and adaptable approach, drawing inspiration from a biological framework that reuses neural network components to build a hierarchy of interacting models, ranging from foundational models to agentic systems (Fig.~\ref{fig:hierarchical_concept}). To quantitatively evaluate its effectiveness, our study concentrates on applying this strategy to science-oriented LLMs, specifically addressing protein mechanics challenges within the broader context of bio-inspired materials analysis and design.

We introduce an algorithm that dynamically combines adapters at the token-level state, allowing for mixing with fine granularity down to individual layers. This technique, termed X-LoRA, provides access to a wide array of functionalities essential for scientific applications of multimodal language models~\cite{Hu2023DeepScience,Buehler2023AMetamaterials,Buehler2022ProteinNetworks,}. 

\subsection{Fundamental concepts of X-LoRA}

We begin by briefly revisiting the principles underlying the development of LoRA adapters. The core idea behind low-rank adaptation (LoRA)~\cite{hu2021lora} posits that weight updates possess a low “intrinsic dimension,” leveraging this by keeping the original weights $W_0 \in \mathbb{R}^{d \times k}$ fixed and limiting the updates to a low-rank decomposition

\begin{equation}
W_0 + \Delta W_0 = W_0 + BA 
\end{equation}

where $B \in \mathbb{R}^{d \times r}$ and $A \in \mathbb{R}^{r \times d}$ with the rank $r \ll {\rm min}(d,k)$. In the training process using LoRA~\cite{hu2021lora}, the original pretrained weights $W_0$ remain fixed, and only the matrices $A$ and $B$ are updated as trainable parameters. The forward pass through a LoRA layer can be represented by the equation:

\begin{equation}
    h = W_0x + \Delta W_0x = W_0x + BAx
\end{equation}

It is important to mention that, in practical settings, LoRA adapters use a fixed scalar weight $\alpha$. This scalar is applied to the decomposed matrices, resulting in

\begin{equation}
    h = W_0x + BAx \cdot \alpha
\end{equation}

Building upon this framework, X-LoRA introduces the concept of employing multiple LoRA adapters, each endowed with distinct capabilities, as demonstrated in prior research on science-oriented LLMs (e.g., \cite{Buehler2023GenerativeDesign}). Rather than statically combining or integrating models, we propose a dynamic gating mechanism that adjusts the scaling of individual LoRA adapters at the level of each token and layer. This approach enables fine-grained mixing deep within the model architecture. Specifically, for a given LoRA adapter $i$, the original scaling parameter $\alpha_i$ is modified by a factor $\lambda_i$, producing a new scaling parameter $\alpha^*_i$ for that adapter:

\begin{equation}
    \alpha^*_i = \alpha_i \cdot \lambda_i
\end{equation}

The scaling factor $\lambda_i$ is generated by a X-LoRA scaling head, which leverages the model’s hidden states and constitutes the trainable element within the X-LoRA framework. Because the scaling head can utilize complex encoded information, it can be efficiently implemented as a straightforward feed-forward neural network with one or a few layers (the X-LoRA implementation allows users to define the specific structure of this neural network).

This scaling computation is performed for each token in the input sequence, providing a high level of granularity by adjusting the scaling of each LoRA adapter operating at various layers in the large language model.

From this point onward, we denote this method of gating the adapter components as scaling. Additional information about this technique is presented in the Materials and Methods section.

\subsection{Paper outline}

The structure of this paper is organized as follows. Initially, we describe the methodology and training framework, including the relevant datasets that lead to the creation of an X-LoRA model with specialized capabilities in the physical sciences, particularly targeting biomaterials such as proteins. Following this, we showcase a range of experiments where the X-LoRA model is applied to diverse tasks, including question answering, conversational and agent-based modeling, protein design and analysis, among others. We provide an in-depth examination of scaling behaviors and confirm the effectiveness of the approach by comparing it with molecular modeling and other physical datasets and techniques.

\section{Results and discussion}

The development of the X-LoRA model involves the following stages:

\begin{enumerate}
    \item Training the foundational base LLM (completed in prior work) \item Separately training a collection of adapters to cultivate expertise across multiple domains (these correspond to distinct or overlapping skill and knowledge areas) \item Training the combined X-LoRA model
\end{enumerate}

Our experimental process begins with training a series of LoRA adapters. We create nine adapters, each fine-tuned with specialized knowledge, based on the Zephyr-7B-$\beta$ model~\cite{HuggingFaceH4/zephyr-7b-betaFace}, which was developed on the foundation of the Mistral-7B model \cite{Jiang2023Mistral7B}:

\begin{enumerate}
    \item Bioinspired materials \item Chain-of-thought (CoT) and reasoning \item Chemistry \item Mathematics \item Physics \item Biology \item Mechanics and materials \item Platypus: Logic and reasoning \item Protein mechanics tasks (including generative sequence-to-property and inverse modeling capabilities)
\end{enumerate}

These focus areas and expertise were selected because our ultimate goal is to develop an LLM centered on STEM with an emphasis on materials science. Nonetheless, we hypothesize that our methodology can be broadly applied across various domains, particularly to establish modalities that facilitate connections between disparate fields of knowledge~\cite{Buehler2023MechGPTModalities,Buehler2023MeLMProblemsb}. 

Following the creation of the LoRA adapters, we proceed to train the X-LoRA model. For this purpose, we utilize a subset drawn from the combined collection of training samples originally employed to train each individual LoRA adapter (for further details, refer to Materials and Methods, which includes comprehensive descriptions of the dataset, implementation, and related information). 

We employ the X-LoRA framework as an experimental system to evaluate its performance in multiple scenarios. We perform a series of experiments, including tackling complex question answering, applying the generative model for protein design, performing analysis and reasoning, and developing more advanced interaction strategies through adversarial agentic modeling. Throughout these analyses, we compare the performance of the base model against the X-LoRA model, monitoring how the scaling head leverages different layers of the various experts during the generation or formulation of responses. 

\subsection{Question answering and observed X-LoRA layer-wise scaling weights}

Our initial investigation focuses on the differences between the X-LoRA model and the underlying foundation model. We pose the following question to both models:

\begin{quote}
\small\texttt{What material behavior (stiffening vs softening) causes supersonic crack speeds in mode II cracks? Think step by step.}
\end{quote}

The base model replies:

\begin{quote}
\small\texttt{Supersonic crack speeds, also known as Rayleigh waves, occur in mode II cracks due to the stiffening material behavior of the material near the crack tip. In mode II fracture, the loading is shear-dominated and the crack propagates perpendicular to the direction of applied stress. Here's a step-by-step explanation:

When the crack tip nears the material, the stresses close to it increase sharply. This leads to an area of elevated stress concentrations, known as the stress singularity.

Within this stress singularity, the material undergoes elastic deformation, resulting in increased stiffness, or stiffening behavior. This stiffening occurs because the material adjacent to the crack tip becomes more restricted as the crack advances [...]}
\end{quote}

Drawing on our own expertise~\cite{Buehler2008AtomisticFailure}, this explanation is not only inaccurate but also unnecessarily verbose. By comparison, the X-LoRA model provides the following response:

\begin{quote}
\small\texttt{The material behavior responsible for supersonic crack velocities in mode II cracks is stiffening.}
\end{quote}

This response is accurate, concise, and clear~\cite{Buehler2008AtomisticFailure}.

The contrast between the base model and the X-LoRA model becomes even more pronounced in specialized protein analysis and design tasks. The prompt given is:

\begin{quote}
\small\texttt{CalculateForceHistory< M S K S N T Y R M L V L L E D D T K I N K E D E K F L K G K P G K M H E F V D E L I L P F N V D E L D E L N T W F D K F D A E I C I P Y V E I L K E S G M K > }
\end{quote}

The base foundation model responds with:

\begin{quote}
\small\texttt{I'm sorry but the input provided is not a valid mathematical expression. It appears to be a sequence of characters that do not form a calculable function. Please provide a valid mathematical expression for me to calculate.

However, I can tell you that the sequence of characters you provided is actually a series of function and variable names that may be part of a larger program or script in a specific programming language. Without context, it's impossible to determine the intended calculation.

If you provide more information, such as the programming language and the context in which this expression appears, I may be able to provide more insight into its meaning and possible calculations. }
\end{quote}

The response from the X-LoRA model is as follows, successfully performing the task of producing the force-deformation curve related to this protein:

\begin{quote}
\small\texttt{[0.013,0.066,0.078,0.045,0.072,0.126,0.127,0.107,0.119, 
0.074,0.106,0.098,0.103,0.128,
0.077, 0.069,0.131,0.108,0.144,0.174, 0.166,0.183,0.185,0.217,0.175,0.230,0.187,0.224, 0.267,0.301,0.315]}
\end{quote}

Figure~\ref{fig:QA_weights_sample_0} presents the results from this question answering alongside the observed X-LoRA scaling weights. The two queries—one centered on mechanics and materials (Fig.~\ref{fig:QA_weights_sample_0}a) and the other on protein mechanics (Fig.~\ref{fig:QA_weights_sample_0}b)—exhibit distinct activation patterns. Notably, we observe that each case employs a complex pattern of scaling values, indicating that the X-LoRA model leverages a heterogeneous combination of adapters distributed across layers. Additionally, the heatmap in Fig.~\ref{fig:QA_weights_sample_0}a demonstrates a wider range of activations spanning numerous experts and layers, including regions of high activation. In contrast, the heatmap in Fig.~\ref{fig:QA_weights_sample_0}b displays a more focused activation pattern, with fewer areas of elevated values. We suggest this might imply that the gating function depicted on the right exercises greater selectivity in activating experts.

In Fig.~\ref{fig:QA_weights_sample_0}, the areas resembling bright lines indicate paths of heightened activation, meaning that the scaling function preferentially activates certain experts at specific layers. The observed patterns imply that importance is not evenly distributed across all experts and layers; rather, it tends to be sparse and selective. This sparsity aligns with findings in other mixtures-of-experts models, where only a subset of experts is chosen for each input to specialize in distinct aspects of the problem space.

By comparing maps across different tasks, both the sparsity and the observed patterns may be crucial for better understanding which experts the model depends on for various input types or at different stages of processing within the neural network (we will later discuss how these maps evolve dynamically during the solving of complex tasks).

To deepen the analysis, we now examine a broader set of questions to assess how the X-LoRA model utilizes the different experts when presented with diverse prompts. We are particularly interested in whether—and how—the X-LoRA model combines various adapters for questions that span multiple skills or domains. Fig.~\ref{fig:QA_weights} presents results from question answering along with the observed X-LoRA scaling weights, corresponding to the queries listed in Table~\ref{tab:table_qa}. In Table~\ref{tab:table_qa}, responses from both X-LoRA and the base model, Zephyr-7B-$\beta$, are shown. The distinctions between the two models are notable.

To clearly illustrate this, we highlight incorrect or questionable responses using \redhighlight{red highlight}. The comparison demonstrates that X-LoRA achieves substantially higher accuracy as well as more concise answers.

As anticipated, we observe intricate mixing of adapters and frequent activation of multiple dominant LoRA experts. For example, Fig.~\ref{fig:QA_weights}d shows that a question about the mechanics of pine cones strongly activates the mechanics/materials adapter alongside the CoT/reasoning and bioinspired adapters. Although X-LoRA selects the correct letter, the answer is not fully accurate since pine cones open under dry, not wet conditions. A follow-up question probing whether the release occurs in dry or wet conditions yields the correct response.

In a protein-related question, Fig.~\ref{fig:QA_weights}f reveals adapter mixing when addressing the mechanics of Bombyx mori silk, activating a combination of the bioinspired, mechanics/materials, and biology adapters. For highly specialized tasks such as protein design (Fig.~\ref{fig:QA_weights}a), we observe predominantly singular activation of one expert.

A more detailed examination of the heatmap of scalings shown in Fig.~\ref{fig:QA_weights}a reveals that besides the clearly prominent importance of the protein mechanics expert, there are several bands where specific experts receive substantial emphasis across multiple layers. This is especially evident for the mechanics/materials and reasoning experts. These bands are not continuous and exhibit a fluctuating pattern, which could indicate particular conditions or features that these experts effectively address. Notably, the reasoning expert (positioned immediately to the left of the protein mechanics expert) is prominent with several layers where it achieves the highest activation, particularly in the upper layers (approximately layers 100 to 140). We suggest that this might point to a crucial role in the model’s final decision-making or output, although further investigation is required to confirm this. Additionally, there is a more localized yet strong activation for the physics and biology experts at certain layers, which might reflect a specialization of these experts for specific features or representations in those regions.

Moreover, Fig.~\ref{fig:QA_token_history} presents a token-level history aggregated over the total scalings across all layers. This analysis displays histories for two examples: Fig.~\ref{fig:QA_token_history}a relates to the task of pine cone seed release, while Fig.~\ref{fig:QA_token_history}b illustrates the outcome of a sequence of protein mechanics tasks. The protein mechanics task consisted of a series of prompt-answer exchanges: initially two property calculations, followed by a generative task aimed at designing a new protein towards the end. Although the usage of experts remains generally consistent, noticeable variations occur in the deployment of the LoRA experts throughout the process. In particular, the protein mechanics expert becomes highly active in the final stages of generation when the new protein is being designed.

Table~\ref{tab:table_qa} demonstrates a variety of use cases spanning different domains and provides a comparison against the performance of the base model alone. We examine several of these cases in greater detail.

For example, we present a question combining protein science, biology, and logic that concerns protein expression patterns. X-LoRA provides the correct answer. In contrast, the base model’s response contains some errors and logical flaws when evaluating the expression of marker proteins in each cell type. The base model accurately recognizes that fibroblasts express Protein B and do not express Protein A due to the mutual exclusivity between Protein A and Protein B. However, it fails to explicitly state that fibroblasts also express Protein C. Since there is no restriction preventing the co-expression of Protein C and Protein B, and fibroblasts have no limitations against Protein C, it logically follows that fibroblasts express Protein C as well. Regarding neurons, the base model correctly begins by stating that neurons express Protein A and not Protein B based on the mutual exclusivity rule. Yet, its discussion about neurons potentially expressing Protein C or both Protein C and Protein B is confusing and flawed. Because neurons express Protein A and not Protein B, according to the rule, the only alternative is for them to express Protein C. Neurons cannot express Protein B, contrary to the initial analysis. Thus, neurons express Protein A and Protein C, but not Protein B. For epithelial cells, the base model correctly observes that they do not express Protein A but do express both Protein B and Protein C. However, the explanation provided is convoluted and includes incorrect claims about the mutual exclusivity between Protein A and Protein C, which was not stated originally. The straightforward logic is that epithelial cells do not express Protein A but must express the two other proteins, Protein B and Protein C, since the original instructions only prohibit co-expression of Protein A and Protein B.

X-LoRA, on the other hand, accurately determines the marker proteins expressed by each cell type by analyzing the given observations and applying logical reasoning. Regarding fibroblasts, X-LoRA correctly infers that fibroblasts do not express Protein A, since the presence of Protein A excludes Protein B expression. Because fibroblasts do express Protein B, and given that each cell type expresses a distinct combination of three marker proteins (although the original statement did not specify the exact number of proteins each cell type expresses, only that they express unique combinations), the conclusion that fibroblasts also express Protein C is justified based on the information provided (aside from the erroneous assumption about the number of proteins, which appears to be a misunderstanding). For neurons, the model accurately identifies their expression of Protein A and Protein C. The rationale is valid, recognizing that Protein B cannot co-express with Protein A, and since neurons express Protein A plus one other marker protein (which is not Protein B), they must express Protein C. Lastly, for epithelial cells, X-LoRA correctly concludes that these cells express Protein B and Protein C. This is because epithelial cells lack Protein A expression, so by elimination, considering they express two marker proteins, those must be Protein B and Protein C.

In the question concerning pine cone release under humid versus dry conditions, the base model produces a confusing and incorrect response. X-LoRA provides a clear and well-reasoned answer. Similar cases exist, such as the question on the exceptional mechanical properties of Bombyx mori silk, where the base model wrongly claims that silk fibers conduct electricity. The base model also struggles with several other domain-specific tasks, for example, the question about the role of hyperelasticity in maximum fracture speeds, where X-LoRA delivers a concise and correct response, while the base model fails to provide an answer. Regarding the strong adhesion of gecko feet, X-LoRA correctly highlights the significance of setae and nanoscale effects, whereas the base model erroneously attributes the adhesion to special-purpose glue proteins. These and other instances discussed in the table and elsewhere in the paper offer compelling evidence for the superior performance of X-LoRA compared to the base model.

Figure~\ref{fig:perf_comp_bioinspiredexam} presents a comparative evaluation of knowledge recall experiments utilizing the bio-inspired knowledge test set introduced in \cite{Luu2023BioinspiredLLM:Materials}. It is evident that X-LoRA achieves the highest performance, despite being a considerably smaller model compared to the BioinspiredLLM or Llama-BioLLM models (7B parameters for X-LoRA versus 13B parameters in the others). Additional comparisons include other recently released LLMs such as Orca-13B and Llama-13b-chat.

The figure also contrasts X-LoRA with the base Zephyr-7B-$\beta$ model, demonstrating that X-LoRA significantly enhances performance. These results suggest that robust performance can be attained using X-LoRA even when smaller base models are employed.

Figure~\ref{fig:image_synth} illustrates an application in image synthesis, where X-LoRA is used to generate a prompt intended for use in other models like DALL-E 3, as shown here.

\begin{quote}
\small\texttt{Create a prompt that I can use for image generation that accurately describes the microstructure of a hierarchical composite.

Explicitly describe the geometry.}
\end{quote}

The comparison (Figure~\ref{fig:image_synth}(b)-(c)) distinctly shows that the prompt generated by X-LoRA is more concise, which leads to a considerably more precise image produced by DALL-E 3~\cite{DALLECard}. Notably, the prompt created by Zephyr-7B-$\beta$ fails to faithfully follow the instruction and introduces various additional design elements like fibers and nanoparticles. These features were absent from the original prompt, which was specifically focused on the microstructure of a hierarchical composite. The X-LoRA prompt is more succinct and targeted, thereby resulting in a significantly improved output.

\subsection{Generative solutions to protein analysis}

We now turn to specific generative and analytical tasks related to proteins. These functions arise in the X-LoRA model through the protein mechanics adapter and are further strengthened by integrating the model’s existing knowledge in biology, bio-inspired materials, reasoning, and logical deduction. This section concentrates on protein-specific tasks alone, providing a quantitative evaluation of this capability.

Fig.~\ref{fig:protein_force_history_prediction} presents outcomes from the protein task aimed at forecasting force-deformation behaviors based on sequences. Fig.~\ref{fig:protein_force_energy_prediction} illustrates the overall accuracy in predicting unfolding force and unfolding energy from the sequence, comparing actual results with model predictions (achieving $R_2=0.85$). In general, the model’s performance is exceptional and, as will be demonstrated in the following section, it also performs very well on the inverse problem of design and cycle-consistent evaluation \cite{Lu2023GenerativePropertiesb}. Notably, unlike previous studies that employed highly specialized GPT-type models for these tasks, our X-LoRA model integrates these specialized functions alongside a range of other deep specializations acquired through various fine-tuned adaptations.

\subsection{Protein design and analysis}

We now broaden the test scenarios to encompass more integrative applications that engage multiple capability domains. Specifically, the experiments described here demonstrate how the X-LoRA model leverages a diverse set of skills and how agentic modeling can generate more complex and comprehensive responses. We begin by addressing a protein design challenge, where the model is tasked with designing a novel protein to achieve a specified target force-deformation behavior. These force-deformation responses have been examined in prior research~\cite{Ni2023ForceGen:Model} and correspond to measurements of force versus deformation as the protein is stretched at both ends, representing a classic protein unfolding experiment~\cite{Lu1998UnfoldingSimulation.}.

Fig.~\ref{fig:design_protein} illustrates the application of the generative protein task. We engineer a protein with a specified force-deformation characteristic (using training data and boundary conditions from the original molecular dynamics (MD) simulations of protein pulling, refer to~\cite{Ni2023ForceGen:Model}), and subsequently evaluate the predicted sequence to determine whether it satisfies the targeted properties (Fig.~\ref{fig:design_protein}a). This cycle-consistent approach is crucial for verifying if the inverse design capability aligns with the forward prediction.

Fig.~\ref{fig:design_protein}b presents the predicted sequence, \texttt{MSKSNTYRMLVLLEDDTKINKEDEKFLKGKPGKMHEFVDELILPFNVDELDELNTWFDKFDAEICIPYVEILKESGMK}. Additionally, Fig.~\ref{fig:design_protein}c displays a folded structure of the protein (generated using AlphaFold 2~\cite{Jumper2021HighlyAlphaFoldb}). To explore the relationship of the designed protein with known sequences, Fig.~\ref{fig:design_protein}d shows a Basic Local Alignment Search Tool (BLAST) Tree~\cite{Altschul1990BasicTool}. Notably, the designed protein exhibits certain connections to hypothetical proteins previously examined in earlier studies but not yet characterized or structurally resolved.

By employing agentic modeling and integrating multiple X-LoRA models as components within a multi-agent framework, we further analyze the predicted sequence using our X-LoRA model. This leverages the advanced capability of our model not only to comprehend protein mechanics but also to synergize this with other knowledge domains. Using the sequence obtained as input, we request an in-depth analysis and suggestions for sequence variations. The task is initiated as follows:

\begin{quote}
\small\texttt{I have identified this amino acid sequence: MSKSNTYRMLVLLEDDTKINKEDEKFLKGKPGKMHEFVDEL ILPFNVDELDELNTWFDKFDAEICIPYVEILKESGMK.

How can I use it?

Provide logical steps, and specifically discuss the engineering aspects of the problem.}
\end{quote}

We conduct multiple agentic interactions to formulate the solution, employing adversarial strategies (see Materials and Methods for details) to encourage the model to probe deeper and address more intricate aspects of the topics. The summary of this dialogue, as produced by the X-LoRA model, is presented below (for the complete conversation, see Supplementary Material, SI Conv Example 1):

\begin{quote}
\small\texttt{The dialogue explores the possible uses of a particular amino acid sequence in the field of protein engineering. This sequence includes both hydrophobic and hydrophilic amino acids, indicating its potential to form a stable conformation and engage with other molecules. To design a protein incorporating this sequence, computational approaches can be employed to forecast the secondary and tertiary structures; however, challenges arise due to the complexity of protein folding and the high computational demands. Enhancements can be achieved by creating more efficient algorithms and integrating new experimental findings. It is essential to validate computational models to guarantee their precision and practical applicability. Methods such as cross-validation, regularization, early stopping, ensemble techniques, and transfer learning are valuable to avoid overfitting and underfitting. Additionally, model interpretation and explanation contribute to guiding practical uses.}
\end{quote}

We also request the model to generate a summary of key points, which leads to:

\begin{quote}
\small\texttt{\begin{itemize}
    \item The amino acid sequence MSKSNTYRMLVLLEDDTKINKEDEKFLKGKPGKMHEFVDELI LPFNVDELDELNTWFDKFDAEICIPYVEILKESGMK holds promise for protein engineering because of its combination of hydrophobic and hydrophilic residues. Computational techniques can predict the protein's secondary and tertiary structures, though challenges include protein folding complexity and computational expense.
    \item Ensuring computational model validation is vital to maintain accuracy and practical value. Approaches like cross-validation, regularization, early stopping, ensemble methods, and transfer learning help prevent both overfitting and underfitting.
    \item Understanding and explaining models further assists in applying them effectively in real-world scenarios.
\end{itemize}
}
\end{quote}

Ultimately, the primary conclusion is:

\begin{quote}
\small\texttt{The key takeaway is that the amino acid sequence MSKSNTYRMLVLLEDDTKINKEDEKFLKGKPGKMHEFVDELILPFNVDELDELNTWFDKFDAEICIPYVEILKESGMK holds promise for protein engineering due to its combination of hydrophobic and hydrophilic residues. To design a protein featuring this sequence, computational approaches can be employed to forecast its secondary and tertiary structures; however, challenges include the complexity of protein folding and high computational demands. It is essential to validate computational models to ensure their accuracy and practical relevance. Methods such as cross-validation, regularization, early stopping, ensemble techniques, and transfer learning are useful for avoiding overfitting and underfitting. Additionally, interpreting and explaining the models can guide practical implementation.}
\end{quote}

We can also delve deeper to examine particular steps in sequence engineering. For this purpose, we request the model to provide targeted design recommendations aimed at enhancing the protein’s stability:

\begin{quote}
\small\texttt{How could I modify the following protein sequence to improve its stability: MSKSNTYRMLVLLEDDTKINKEDEKFLKGKPGKMHEFVDELILPFNVDELDELNTWFDKFDAEICIPYVEILKESGMK? }
\end{quote}

A summary of the discussion is as follows (for the complete exchange, refer to Supplementary Material, SI Conv Example 2):

\begin{quote}
\small\texttt{The discussion centers on strategies to increase protein sequence stability by adding hydrophobic amino acids and disulfide bonds. The biologist notes that hydrophobic residues generally enhance protein stability, while disulfide bonds also contribute positively. The engineer requests concrete examples of hydrophobic amino acids and guidance on the number of disulfide bonds to incorporate. The biologist cites examples of proteins stabilized through such modifications, including insulin, bovine pancreatic trypsin inhibitor, and green fluorescent protein. The conversation further addresses potential challenges associated with adding hydrophobic residues or disulfide bonds, such as alterations in protein function, increased structural complexity, and risk of aggregation. The biologist suggests these issues can be mitigated by judicious amino acid selection, computational modeling, and experimental validation.}
\end{quote}

The summary of key points:

\begin{quote}
\small\texttt{\begin{itemize}
\item Enhancing protein sequence stability can be achieved by incorporating hydrophobic amino acids and disulfide bonds.

Hydrophobic amino acids generally increase protein stability, and disulfide bonds contribute similarly.

Proteins like insulin, bovine pancreatic trypsin inhibitor, and green fluorescent protein have been successfully stabilized by applying these modifications.
\item Possible downsides of adding hydrophobic amino acids or disulfide bonds include altered protein function, greater complexity, and a risk of aggregation.
\item These issues can be mitigated through careful amino acid selection, computational modeling, and experimental testing.
\end{itemize}
}
\end{quote}

The main conclusion is:

\begin{quote}
\small\texttt{The key point from the discussion is that to stabilize the protein sequence MSKSNTYRMLVLLEDDTKINKEDEKFLKGKPGKMHEFVDELILPFNVDELDELNTWFDKFDAEICIPYVEILKESGMK, one can introduce hydrophobic amino acids and disulfide bonds. This directly addresses the initial inquiry.}
\end{quote}

From the standpoint of protein engineering, these recommendations are logical as they would likely result in more stable protein designs. Concrete examples and methodologies are developed and suggested for further exploration by researchers.

In another scenario, we analyze three sequences. We first ask the model to assess the strength of each sequence, then to reason about the outcomes: identify the strongest candidate and explain why this protein likely has the highest strength. Additionally, we request the model to describe the steps needed to synthesize it in the laboratory. The three sequences (manually derived by modifying one of the test set sequences) are:

\begin{quote}
\small\texttt{
\begin{itemize}
    \item LNTWFDKFDAEICIPYVEILKESGMK \item AITWFDKFDAEICIPYVEALKIAAMK \item AAAAAAAAKFDAAEICIIAAAAAAAAKIAAMK
\end{itemize}
}
\end{quote}

This example evaluates whether the model can integrate multiple skills and apply them flexibly as the conversation unfolds. Specifically, it tests the ability to switch dynamically among capabilities such as understanding protein mechanics, conducting rational analysis, generating new hypotheses, and so forth.

The conversation is summarized in Figure~\ref{fig:conversation_evolve}, which also illustrates how the scalings evolve throughout the discussion (the analysis includes both the deep layer-wise adapter scaling and an aggregated assessment of these to provide an indicator of the effective utilization of the different experts). The findings clearly demonstrate that the protein task is initially highly relevant, progressively transitioning to a more dominant presence of the bioinspired adapter and the biology adapter. By the conclusion of the conversation, the biology adapter emerges as the most dominant.

To evaluate the predictions and assess the response quality, the three proteins involved in this task are shown in Figure~\ref{fig:protein_visuals} (proteins folded using Alpha Fold 2~\cite{Jumper2021HighlyAlphaFoldb}, via ColabFold \cite{Mirdita2022ColabFold:All}). We observe that the most stable structure in the center displays the most well-organized secondary structure, consistent with the idea that it exhibits the highest unfolding force, as predicted by the X-LoRA model in Figure~\ref{fig:conversation_evolve}. The other two proteins are less structured; the first is predominantly helical with some turns, and the last is partially unstructured with a bend in its geometry. These characteristics account for the lowest unfolding force predicted by the model.

By examining the visual depictions of the folded structures, we verify that the most stable structure in the center exhibits the most ordered secondary organization, consistent with the idea that it produces the highest unfolding force as predicted by the model. We observe that the other two models are less structured. The first one is predominantly helical with some turns, while the last one is partially disordered and features a kink in its geometry. We believe these characteristics likely account for the lowest unfolding force observed, aligning with the predictions of the X-LoRA model.

In a second example, we focus on the energetic characteristics of proteins, specifically the unfolding energy as the protein is stretched from its initial conformation to a fully unfolded state~\cite{Ni2023ForceGen:Model}. We again analyze three sequences (the first two are the same as previously considered, whereas the third sequence is longer, representing a combination of the second sequence at the start and end and the first sequence in the middle):

\begin{quote}
\small\texttt{\begin{itemize}
    \item LNTWFDKFDAEICIPYVEILKESGMK \item AITWFDKFDAEICIPYVEALKIAAMK \item AITWFDKFDAEICIPYVEALKIAAMKLNTWFDKFDAEICIPYVEILKESGMKAITWFDKFDAEICIPYVEALKIAAMK
\end{itemize}
}
\end{quote}

We initially instruct the model to calculate the unfolding energy for each sequence, then to analyze the outcomes: based on the amino acid sequence, we request an explanation for why a particular protein is probably the most stable. Finally, we query the model about the potential functions of the most stable protein.

Figure~\ref{fig:MJB_20} presents a transcript of a discussion that combines various specialized experts and synthesizes knowledge across these domains, demonstrated here in the context of analyzing the unfolding energy of three proteins. By the conclusion of the dialogue, the CoT/reasoning adapter becomes the most dominant component as the X-LoRA model answers a question requiring reasoning over the results to predict probable protein structural characteristics and possible functions. Specifically, the mechanistic inquiry regarding the source of protein stability highlights the formation of hydrophobic interactions among nonpolar amino acids and hydrogen bonds between polar amino acids. The model forecasts that several hydrophobic amino acids (I, F, W, Y) and hydrogen-bonding amino acids (D, E, K) play roles in enhancing protein stability.

Additionally, the model anticipates that the inclusion of cysteine (C) residues in the sequence promotes the formation of disulfide bonds, which further reinforce the protein’s structural stability. Notably, the key structural feature predictions, which explain the protein's high stability, are accurately identified, as corroborated by structural analyses displayed in Fig.~\ref{fig:MJB_21}. This analysis confirms the presence of all these features. It also clarifies why the other two proteins, composed of considerably shorter segments with varying helical domain stability, exhibit substantially lower unfolding energies.

\subsection{Adversarial agentic modeling to connect distinct scholarly disciplines and knowledge yielding ontological knowledge graph generation}

Given that our X-LoRA model possesses extensive expertise across multiple fields, it can be employed to explore links among disparate concepts, knowledge bases, and specialized domains. In one experiment, the model is queried twice with similarly structured questions but each emphasizing a different perspective. Subsequently, we generate triplets to produce an ontological knowledge graph that transforms the answers into a more organized format~\cite{Giesa2012CategoryDesign,Spivak2011ReoccurringAnalogies}. This resulting graph delivers a cohesive understanding of the insights generated and visually represents the relationships between concepts in an interpretable and mechanistic way.

To address the question, we employ two adversarial agents. Agent 1 is characterized as a "critical physicist," while Agent 2 assumes the role of a "philosopher." Agent 1 initiates the inquiry and is responsible for probing and generating further questions in response to every answer provided by Agent 2. As the dialogue progresses, it delves deeper into the topic, revealing multiple facets. The question is deliberately crafted to span various disciplinary domains to evaluate the model's capability to synthesize diverse ideas and concepts.

The question is:

\begin{quote}
\small\texttt{Consider an abstract representation of a classical musical composition as a graph of interacting notes, chords, melodies, rhythms and so on.

What happens if we would apply pressure, up to a point of failure, on this abstract concept of music?

Provide logical steps, mathematical reasoning, and discuss specifically the physics of the problem.}
\end{quote}

The complete conversation is provided in the Supplementary Material (refer to SI Conv Example 3). For brevity, we present here only the summary and main takeaway, capturing the key elements:

\begin{quote}
\small\texttt{The discussion explores the effect of applying pressure to an abstract model of a classical musical composition, viewed as a physical system composed of interacting notes, chords, melodies, and rhythms. Pressure is represented quantitatively through stress, with the system's failure point identified when the stress attains a critical threshold. To validate the mathematical model's predictions, a physical analog of the abstract musical composition can be constructed, enabling experimental observations to be compared against theoretical forecasts. Limitations of the proposed model include its simplified depiction of interactions, static nature, assumption of homogeneous distribution, and linear correlation between stress and pressure. Enhancing the model's precision and reliability would require incorporating more detailed and realistic interactions, dynamic behavior, heterogeneous distributions, and non-linear relationships.

[...]

The most significant conclusion is that exerting pressure on an abstract representation of a classical musical composition can be conceptualized as modifying the force field among the interacting notes, chords, melodies, and rhythms. Increasing pressure raises the system's internal stress, which, upon reaching a critical threshold, causes the system to fail and collapse under its own structural constraints. This explanation addresses the original question by offering a physical rationale for how pressure influences the abstract notion of music and how this influence can be quantified through stress.}
\end{quote}

Although the generated text yields valuable insights, constructing an ontological representation of the data would facilitate more interpretable analysis~\cite{Buehler2023MechGPTModalities}. For graph creation, we merge the entire dialogues from both agents to build a knowledge graph. An example of such a graph is displayed in Figure~\ref{fig:graph}. The resulting graph is partitioned into communities using the Girvan-Newman algorithm~\cite{Girvan2002CommunityNetworks}, identifying 12 distinct communities in total. The figure presents both the comprehensive graph (Figure~\ref{fig:graph}a) and detailed views with labeled nodes (Figures~\ref{fig:graph}b-c).

\subsection{Creation of X-LoRA-Gemma Incorporating Combined Protein, Chemical, Bio-Inspired, and Materials Mechanics Functionalities}
\label{gemma-section}

To demonstrate the applicability of the proposed methodology to alternative base models with different architectures, we trained an additional X-LoRA model, this time using the Gemma-7B-it model~\cite{Google/gemma-7b-itFace} as the foundation. In this instance, we utilized a collection of four LoRA adapters, defined as follows:

\begin{enumerate}
    \item Bioinspired materials \item Mechanics and materials \item Protein mechanics tasks (including generative sequence-to-property and inverse design functionalities) \item Quantum mechanics-based molecular properties QM9 (featuring generative SMILES-to-property and inverse design capabilities; refer to Table~\ref{tab:table_QM9} for a definition and summary of all properties computed in that dataset)~\cite{Ruddigkeit2012EnumerationGDB-17,Ramakrishnan2015ElectronicSpace}
\end{enumerate}

The X-LoRA-Gemma model was developed in a manner similar to the initial X-LoRA model but seeks to incorporate distinct capability sets, particularly by adding the function of predicting 12 quantum mechanical properties and designing molecules to meet those 12 quantum mechanical property criteria. This new adapter is trained concurrently with the other adapters mentioned above, which are based on the same datasets employed in the earlier experiment. It is important to note that while the majority of this paper centers on the Zephyr-based X-LoRA model, within the context of this section, we differentiate the models by referring to the original as X-LoRA-Zephyr and the Gemma-based model as X-LoRA-Gemma.

Figure~\ref{fig:protein_force_energy_prediction} displays the overall performance of X-LoRA, illustrating predictions of unfolding force and unfolding energy derived from the sequence ($R_2 = 0.87$), alongside predictions of quantum mechanical molecular properties from the QM9 dataset ($R_2 = 0.96$). We observe outstanding performance on both tasks, with slightly improved outcomes for the protein design task compared to the original X-LoRA model.

Next, we showcase the capabilities of X-LoRA-Gemma in a straightforward multi-agent design scenario. We examine a small organic molecule, \texttt{CCCC12CC1C(=O)N2} (selected from the QM9 dataset test set). This bicyclic molecule contains a ketone (=O) and an amine (N) integrated into the ring structure (see Figure~\ref{fig:QM9_design}(a) for a visualization). We first pose the following prompt to the X-LoRA-Gemma model:

\begin{quote}
\small\texttt{Here is a molecule in SMILES code: CCCC12CC1C(=O)N2

What modifications would be necessary to make this molecule more suitable for polymer design?}
\end{quote}

The response is:

\begin{quote}
\small\texttt{To enhance the utility of this molecule for polymer design, you might:
\begin{enumerate}
\item  Add a reactive functional group, like an amine or hydroxyl group, to enable further modifications and crosslinking.
\item  Extend the carbon backbone to be longer and more flexible, improving the molecule's capacity to form polymers with diverse architectures and characteristics.
\item  Attach a side chain containing a functional group, such as a carboxyl group, to facilitate covalent bonding with other molecules and polymers.
\item  Introduce a cyclic structure, for instance a lactam ring, to produce a polymer that is more rigid and stable.
\item  Incorporate various side chains to generate a polymer with a broad spectrum of properties and functionalities.
\item  Vary the number of carbons in the backbone to enable the synthesis of polymers with differing molecular weights and attributes.
\item  Insert a crosslinking site within the molecule to create a polymer that is more stable and less reactive.
\item  Include a repeating unit to develop a polymer with defined molecular weight and properties.
\item  Combine several of these strategies to design a versatile and practical polymer suitable for multiple applications.
\end{enumerate}}
\end{quote}

Next, we concentrate on the initial recommendation.

\begin{quote}
\small\texttt{What approaches can I use to enhance the overall reactivity of a molecule? I intend to utilize this molecule for constructing a polymer material.

The molecule of interest, represented in SMILES notation, is: CCCC12CC1C(=O)N2

I have the capability to adjust the following properties:

Dipole moment, mu Polarizability, alpha Highest Occupied Molecular Orbital energy, homo Lowest Unoccupied Molecular Orbital energy, lumo HOMO-LUMO gap: The energy difference between HOMO and LUMO, gap Electronic spatial extent Zero-point vibrational energy Heat capacity at constant volume Internal energy at 0 K Internal energy at 298.15 K Enthalpy at 298.15 K Free energy at 298.15 K Which are the top THREE properties I should modify, and should I increase or decrease them? }
\end{quote}

The feedback from X-LoRA-Gemma is:

\begin{quote}
\small\texttt{\begin{itemize}
    \item  HOMO-LUMO gap (reduction) \item  Dipole moment (enhancement) \item  Polarizability (enhancement)
\end{itemize}
Altering these properties can enhance the overall reactivity of the molecule, making it more appropriate for use as a polymer material building block. }
\end{quote}

We leverage the generative and predictive functionalities of X-LoRA-Gemma to design molecules exhibiting these desired property changes.

For the optimization, we employ a two-agent system composed of two instances of X-LoRA-Gemma. One agent performs the generative task \texttt{GenerateMolecularProperties<...>} to create candidate molecules. Subsequently, the other agent calculates properties using \texttt{CalculateMolecularProperties<...>} and evaluates the discrepancy between the target and predicted properties. This cycle continues until either the maximum iteration count is reached or the mean-squared error falls below a predefined threshold.

The outcomes of the full procedure are illustrated in Figure~\ref{fig:QM9_design}. Panel (a) presents the target design alongside its deviation from the original molecule's properties. Panel (b) displays the top 20 generated designs, accompanied by the mean-squared error distribution in panel (c) and the mean-squared error progression over designs in panel (d). The optimal design, \texttt{CC(C)C1=CC(=O)NC=C1}, aligns closely with the target, as demonstrated in panel (e). Lastly, panel (f) depicts an optimized 3D molecular structure, obtained using UFF~\cite{Rappe1992UFFSimulations} to refine the geometry.

To evaluate the molecule's reactivity, we also calculate the Gasteiger charges~\cite{Gasteiger1980IterativeCharges}, with a visual summary in Figure~\ref{fig:QM9_design}(f). The molecule includes an aromatic ring featuring a ketone group (C=O), which is more electrophilic than typical carbon-carbon or carbon-hydrogen bonds, rendering it prone to nucleophilic attacks. The combination of the nitrogen-containing aromatic ring and the ketone forms a pyridone ring structure. This ring exhibits greater reactivity than a simple aromatic ring because the heteroatom (nitrogen) and ketone introduce reactive sites, such as susceptibility to nucleophilic addition at the carbonyl carbon and electrophilic substitution on the ring. The nitrogen atom, as part of the heteroaromatic system, influences aromaticity and reactivity by modulating electron density and distribution within the ring. This alters how the molecule interacts with other species, making the nitrogen-containing ring a probable locus for chemical reactions, especially electrophilic attacks facilitated by the electron-donating character of the nitrogen.

The carbon atom within the ring, situated next to both an oxygen and a nitrogen atom, carries a partial positive charge (+0.25), indicating it is somewhat deficient in electron density. This arises because oxygen, being more electronegative, draws electron density away from the carbon, rendering it less electron-rich. This effect is partially offset by the neighboring nitrogen, which is less electronegative than oxygen but more so than carbon. Nevertheless, the partial positive charge on the carbon suggests it is a probable site for nucleophilic attack, where a nucleophile (an electron-rich species) would be attracted to this electron-poor carbon.

The nitrogen atom, possessing a partial negative charge (-0.33), signifies it is relatively electron-rich compared to its typical state. This can be attributed to its lone pair of electrons and the influence of the aromatic ring's electron system. In aromatic heterocycles, nitrogen’s lone pair can engage in the delocalized $\pi$-electron network, although the precise involvement depends on the ring’s structure and the electronegativities of neighboring atoms. The partial negative charge implies that nitrogen is more nucleophilic, meaning it has an increased tendency to donate electrons, either by bonding with electrophiles or participating in protonation reactions. Additional quantum mechanical analyses are needed to explore this further.

The carbon atom bearing the partial positive charge represents a reactive center for nucleophilic attacks. Other molecules may interact at this site via mechanisms where nucleophiles target electron-deficient carbons.

The hydrogen atom attached to the nitrogen, which has a partial positive charge of 0.17, is a favorable candidate for forming hydrogen bonds with other atoms or molecules. Hydrogen bonds are a type of dipole-dipole interaction that occurs when a hydrogen atom bonded to a highly electronegative element (such as nitrogen, oxygen, or fluorine) interacts with another electronegative atom possessing a lone pair of electrons. This attribute enables the molecule to act as a hydrogen bond donor, a key factor influencing its solubility in water and other polar solvents, as well as its behavior in biological environments.

The ketone group positioned next to the nitrogen-containing aromatic ring imparts distinctive reactivity that can be harnessed in the synthesis of polymers~\cite{McQuarrieSimonPhysicalChemistry,CareySundbergAdvancedOrganicChemistry,Varghese2022BeyondApplications,Varghese2022BeyondApplications}. This compound belongs to the lactam class, characterized by a cyclic amide formed between the ketone and nitrogen within the ring. Lactams play a significant role in polymer chemistry, especially in producing polyamides, a polymer class widely utilized in various material applications. The ketone functionality can participate in multiple chemical reactions beneficial for polymer chemistry, including nucleophilic addition. Such reactions could facilitate polymer chain growth or introduce side chains that alter the polymer’s characteristics. Additionally, the nitrogen atom within the aromatic ring can markedly influence the polymer’s electronic properties and engage in hydrogen bonding, affecting attributes like melting point, solubility, and mechanical behavior. N-heteroaromatic compounds are also recognized for their capacity to coordinate metal ions, enabling potential uses in catalysis or the development of materials with electronic or photonic functionalities. Furthermore, this molecule may undergo polymerization through its amide (lactam) group; polyamides derived from such polymerizations exhibit notable strength, elasticity, and resistance to wear and chemicals, finding extensive use in diverse applications ranging from textiles (e.g., nylon) to automotive components~\cite{Varghese2022BeyondApplications,Varghese2022BeyondApplications}.

Alkyl substituents can enhance the hydrophobic nature of the molecule, leading us to predict that polymers synthesized from monomers containing these alkyl groups will exhibit increased hydrophobicity. This alteration influences their solubility in various solvents and their interactions with water, which may be beneficial for applications that demand water resistance or particular solubility profiles. Additionally, alkyl groups can influence the polymer’s mechanical properties. For example, they may serve as plasticizers within the polymer matrix, thereby enhancing the polymer’s flexibility. This effect could produce materials that are more flexible or possess improved impact resistance, depending on the alkyl groups’ chain length and branching.

\subsubsection{Additional benchmarking} We present further benchmarking results for the X-LoRA-Gemma model, illustrated in Figure~\ref{fig:perf_MMLU}, using a selected subset of the MMLU benchmark focused on chemistry, physics, and biology (domains closely related to those on which our model was trained)~\cite{Hendrycks2020MeasuringUnderstanding}. Our findings show that both X-LoRA models analyzed in this study outperform their respective base models. Specifically, the X-LoRA models surpass the base models’ performance (Zephyr model: 54.6\% for the base versus 56.6\% for X-LoRA, Gemma model: 40.6\% for the base compared to 52.4\% for X-LoRA-Gemma). The highest overall accuracy is achieved by the X-LoRA-Zephyr model, though both models exhibit comparable results.

\section{Conclusion}

Multimodal LLMs have numerous current and prospective applications in scientific fields, coupled with a strong demand for highly specialized, task-specific expertise. Beyond general language processing tasks, the capability to perform computational or generative functions, as illustrated here with protein mechanics, is essential~\cite{Luu2023BioinspiredLLM:Materials,Buehler2023MechGPTModalities,Buehler2023MeLMProblemsb}. X-LoRA, which is based on open-source language models, allows for the dynamic blending of expertise and knowledge across different domains, as shown in several examples including Figure~\ref{fig:conversation_evolve}. This study demonstrated how the X-LoRA model autonomously and adaptively evolves its use of adapters, mixing particular adapters in a sophisticated manner throughout the LoRA layers. This intricate mixing of LoRA layers is consistently observed across all examples analyzed here, such as in Figures~\ref{fig:QA_weights_sample_0} and \ref{fig:QA_weights}. This phenomenon presents an intriguing challenge to better understand the foundational mechanisms from a physics standpoint, conceptualizing a LLM as a large-scale graph-formation model exploiting complex graph structures. Consequently, the blending of adaptation experts at deep layer levels, as performed here, modifies these mechanisms in a way that promotes the development of effective solutions.

The proposed algorithm enhances the traditional inference approach of a single forward pass by incorporating two forward passes. This trains the model to first "reflect" on the question and how it might reconfigure itself before providing an answer. This approach realizes a simple form of "self-awareness," enabling the model to adjust its own architecture to optimally address a given task. As a result, despite having a relatively modest parameter count of 7 billion, the model is capable of reasoning across a wide range of scientific domains (including biological materials, mathematics, physics, chemistry, logic, mechanics, and more), greatly improving its ability to produce innovative solutions. Furthermore, it could inspire the design of alternative model architectures, such as approaches that extend the reconfiguration strategy to parts or the entirety of non-adapted systems. In these cases, the scaling head would dynamically reorganize sections of a trained model or integrate multiple models. As demonstrated with X-LoRA, such strategies serve as powerful tools to broaden or combine the capabilities of different models.

We conducted multiple performance evaluations, including a comparison with significantly larger models (Figure~\ref{fig:perf_comp_bioinspiredexam}), where X-LoRA demonstrated a clear superior performance. The comprehensive comparison presented in Table~\ref{tab:table_qa} revealed that X-LoRA substantially outperforms the base model across a diverse range of tasks and application areas. Additional comparisons involved quantitative evaluations of numerical prediction tasks, such as protein mechanics and the quantum-mechanical properties of molecules, where X-LoRA achieved high accuracy with $R_2$ values ranging from 0.85 to 0.96. For reference, these results exceeded other published findings, including protein unfolding force prediction using the MeLM model~\cite{Buehler2023MeLMProblemsb}, which reported an $R_2=0.78$. The ability to accurately address multiple prediction tasks, as well as inverse design and advanced reasoning capabilities, was highlighted in the molecular design example (Section~\ref{gemma-section}). This example also demonstrated that X-LoRA can be effectively implemented using various base model architectures (Gemma versus Zephyr/Mistral). Future research could extend this approach further by developing X-LoRA models for even larger base models.

One limitation of the proposed method is the necessity for two forward passes, which may result in increased computational cost: the first to compute hidden states that feed into the X-LoRA scaling head, and the second with $\Lambda$ applied to the adapters. However, the initial pass serves the X-LoRA scaling head and leverages fundamental LLM capabilities, unlike approaches that require learning knowledge separately by training a larger scaling head. For practical deployment with model-serving systems, distinct key-value caches need to be maintained for both the scaling and forward passes. To enhance inference speed, we created \texttt{Mistral.rs}, an LLM serving platform implemented in Rust~\cite{matsakis2014rust} that supports X-LoRA and incorporates multiple performance optimization techniques. Nevertheless, a significant advantage of our method is its compatibility with any model without needing to modify internal logic. We consider this a benefit, especially given its applicability to the extensive resources available within the Hugging Face ecosystem.

One benefit of the proposed method is that it offers a straightforward way to integrate with any existing LLM (for example, since the X-LoRA code has been made available, it should be readily applicable to any model within the Hugging Face framework). Another benefit lies in the scaling head's ability to leverage the inherent strengths of the LLM and learn sophisticated strategies for optimally combining adapters to produce specific answers. When training X-LoRA, it is advisable to use complex examples of question-answer pairs or conversations that enable the model to discover the most effective combinatorial approach. We conducted several analyses focusing on protein design and mechanics. For example, Figure~\ref{fig:MJB_20} and Fig.~\ref{fig:MJB_21} showcase a stability analysis of multiple protein sequences, including comparisons, reasoning about the predicted outcomes, and a mechanistic explanation that offers a bottom-up chemical basis for the predicted behavior. As demonstrated in Fig.~\ref{fig:MJB_21}, this was directly validated through a structural analysis of the protein’s folded 3D conformation.

Regarding future directions, further studies could investigate a broader range of adapter experts, expanding upon the initial expertise encompassed in our LoRA set. The X-LoRA model surpasses the performance of the base model alone or the base model with a single adapter, delivering answers to complex questions by drawing on the specialized capabilities provided by the set of adapters. There also appear to be synergistic effects among the different adapters, as suggested by the intricate mixing outcomes of scaling weights across layers. This aspect warrants additional exploration and could be compared to approaches such as SLERP, which combines models using different techniques. While we trained the X-LoRA scaling head using a subset of the training data—consisting of several hundred samples from each of the original training datasets used to develop the adapters—more targeted creation of a training dataset designed for this purpose could be established. For instance, we observed that multiple-choice questions tend to activate the reasoning-focused adapters, since these adapters were trained on such data. Additional prompting is necessary to assist the scaling head in grasping a specific domain relevant to the questions posed, potentially through particular system messages or alternative strategies. More broadly, the creation of suitable training datasets is a promising area for further research, including methods tailored to effectively induce mixing of layer-wise scaling. Our experiments tracking these mechanisms over extended conversations indicate that the model can dynamically adjust scaling strategies to optimally address certain tasks, which is an encouraging outcome.

We identified intriguing activation patterns of experts across different layers, with compelling insights elaborated in the paper. This analysis could be extended further and might provide valuable understanding of the benefits and mechanisms behind mixing components or groups of layers in models. Such exploration would be especially compelling when applying the X-LoRA technique to fields beyond protein mechanics and protein design, as demonstrated here. We leave this exploration for future work.

Another significant direction involves deeper investigation into agentic modeling, exemplified here by employing two adversarially interacting X-LoRA models. Future research could incorporate more than two models to enrich their interactions and introduce new functionalities, thereby driving models into novel generative domains. This approach may also facilitate developing methods that combine physics-based modeling or other validation processes, potentially leveraging code-writing/executing agents or agents employing function calling or other processing strategies~\cite{Ghafarollahi2024ProtAgents:Learning}.

\section{Materials and Methods}

The X-LoRA code and example implementations are accessible at \url{https://github.com/EricLBuehler/xlora}. Model weights and datasets linked to the X-LoRA presented in this study, along with additional examples, can be found at \url{https://huggingface.co/lamm-mit/x-lora}. The X-LoRA-Gemma model is available at \url{https://huggingface.co/lamm-mit/x-lora-gemma-7b}. The \texttt{Mistral.rs} inference engine is provided at \url{https://github.com/EricLBuehler/mistral.rs}.

\subsection{Mathematical formulation of X-LoRA} We define the $\mathbf{scalings}$ as a matrix comprising token- and layer-level coefficients for the LoRA adapters. The X-LoRA $\mathbf{scaling}$ $\mathbf{head}$ processes the base model's hidden states by applying fully connected linear layers to predict the scalings, followed by a softmax function across the final dimension to ensure that scalings over all experts sum to one. ReLU activation functions and dropout layers are interspersed between these linear layers in the X-LoRA scaling head. The hidden layer size can be configured to effectively reduce the dimension from the hidden state size to the scaling dimension, or to adopt a widening-then-narrowing architecture akin to the feed-forward layer in transformer blocks.

The scalings form a tensor $\Lambda \in \mathbb{R}^{s \times l \times n}$, where $s$, $l$, and $n$ denote the sequence length, number of layers, and number of adapters, respectively. For each layer, the relevant scalings are extracted as a matrix $\lambda \in \mathbb{R}^{s \times n}$. These scalings $\lambda$ are applied to each LoRA adapter by broadcasting the associated X-LoRA scaling $\lambda_i \in \mathbb{R}^{1 \times 1 \times s}$ to the inputs of the decomposition matrices $A$ and $B$.

Given $W_0 \in \mathbb{R}^{d \times k}$ and adapter-specific matrices $B_i \in \mathbb{R}^{d \times r_i}$ and $A_i \in \mathbb{R}^{r_i \times d}$ with rank $r_i \ll {\rm min}(d,k)$, the forward operation of an X-LoRA layer is described by:

\begin{equation}
    h = W_0 x + \sum_{i=1}^{n} B_i A_i (x \cdot \lambda_i) \alpha_i
\end{equation}

Examples of these scaling values are illustrated in the main manuscript, visualized either through bar plots or using the \texttt{contourf} function from Matplotlib~\cite{MatplotlibDocumentation}.

\subsection{X-LoRA implementation algorithm and code} The implementation presented here prioritizes user-friendliness and flexibility, enabling straightforward adaptation to various models, particularly those within the Hugging Face ecosystem. The X-LoRA method functions by:

\begin{enumerate}
    \item Keeping the base model and adapters fixed while tuning for improved performance.
    \item Combining LoRA adapters based on the scaling predictions provided by the X-LoRA scaling head.
\end{enumerate}

All code is developed in Python using PyTorch~\cite{Paszke2019PyTorchLibrary}. 

\subsubsection{Dual forward pass strategy for self-aware inference} Since X-LoRA necessitates computing hidden states to forecast the scaling factors, we adopt a dual forward pass technique. The first pass allows the model to analyze the task and decide how to adjust itself, followed by the actual inference step. This approach introduces a form of `self-awareness' that enables X-LoRA to allocate additional computation time for internal reflection rather than immediate output generation.

We refer to the initial pass as the scaling pass and the subsequent one as the forward pass. During the scaling pass, the scalings $\Lambda$ are set to zero. We also explored setting $\Lambda$ to ${n}^{-1}$ and observed similar outcomes, indicating stable performance (this parameter can be configured in the X-LoRA package). To facilitate application across different models, we implement a forward pass hook. This mechanism is crucial in linking the scaling and forward passes, and it is the principal feature that supports compatibility with all models.

\begin{enumerate}
    \item $\mathbf{Scaling}$ $\mathbf{pass}$: Executed by running the base model with adapters fixed at constant values. The resulting hidden states feed into the X-LoRA scaling head to produce $\Lambda$.
    \item $\mathbf{Forward}$ $\mathbf{pass}$: Utilizes $\Lambda$ to blend the adapters during the model's forward pass. The generated logits are then outputted as the result of the X-LoRA model.
\end{enumerate}

It is important to note that sharing the same key-value cache between the scaling pass and the forward pass may disrupt proper inputs to the scaling head, as the cache persists and accumulates over both passes. Consequently, we avoid using the key-value cache during training and inference, implementing appropriate management of this in \texttt{Mistral.rs}.

\subsubsection{Freezing weights and selectively enabling model segments for training} In an X-LoRA model, the base model and adapters remain frozen (meaning their weights are not updated during training), making the X-LoRA scaling head the sole trainable component. The provided code also allows the option to unfreeze all adapters along with the X-LoRA scaling head for training. Furthermore, one might choose to train the entire model simultaneously—including the X-LoRA scaling head, adapters, and the underlying foundation model—potentially applying distinct learning rates to each. Another avenue for future research involves extending the concepts introduced here to devise scaling functions connecting two or more foundation models, enabling a traceable form of SLERP-style modeling. This prospect remains to be explored in subsequent studies.

\subsubsection{X-LoRA layers} An X-LoRA layer is designed to scale and combine outputs from multiple adapters. It encapsulates the original LoRA layer, thereby maintaining access to the specific adapters associated with that layer.

To apply X-LoRA seamlessly to any model, the algorithm replaces the original LoRA layers by generating X-LoRA layers and subsequently injects their forward pass logic into the LoRA layer. This process utilizes Python’s object model capabilities to facilitate the layer swapping mechanism.

In summary, during the forward pass, each X-LoRA layer performs the following operations (refer to Figure \ref{fig:general_arch}):

\begin{enumerate}
    \item $\mathbf{Scaling}$: Multiply each LoRA adapter’s input signal by its corresponding scaling factor $\lambda_i$.
    \item $\mathbf{Forward}$: Execute the forward pass of the wrapped LoRA layer.
\end{enumerate}

\subsection{Efficient Inference with \texttt{Mistral.rs} implemented in Rust} To enhance inference efficiency, we developed \texttt{Mistral.rs}, a large language model serving platform implemented in Rust~\cite{matsakis2014rust} that supports X-LoRA and incorporates multiple performance optimizations. Among the implemented methods, without implying any priority, are:

\begin{itemize}
    \item LoRA adapter weight stacking: When the matrices $A_i$ and $B_i$ of all LoRA adapters share the same dimensions, we can concatenate these LoRA adapter weight matrices $A_i$ and $B_i$ along their first dimension.

Mathematically, given $W_0 \in \mathbb{R}^{d \times k}$, for adapter matrices $A_i \in \mathbb{R}^{r \times d}$ and $B_i \in \mathbb{R}^{d \times r}$ with $r_i \ll {\rm min} (d,k)$, we concatenate these matrices so that $A\in \mathbb{R}^{i \times r \times d}$ and $B\in \mathbb{R}^{i \times d \times r}$. The parameters $\alpha_i$ are also applied to the stacked matrices during initialization.

To apply scalings, we rearrange the dimensions so that $\lambda \in \mathbb{R}^{i \times 1 \times 1}$. These are then applied via broadcast multiplication.

\item Non-granular scalings To minimize how often the scaling pass runs, we introduce an option for non-granular scalings. This approach caches the scalings after $n$ completed runs because the KV cache guarantees that the sequence length remains one for the rest of the request. Utilizing this method significantly lowers latency by reducing the number of calls to the base model.

\item Fused Compute Unified Device Architecture (CUDA) kernels To enhance the forward pass efficiency, we utilize fused CUDA kernels. Such fused kernels allow combining complex operations, like Rotary Embedding (\cite{su2021roformer}), which would otherwise be executed sequentially, into a single kernel that can run in parallel on a GPU. Our kernels are based on the vLLM implementation (\cite{kwon2023efficient}).

\item Quantization The inference engine supports quantization, enabling the use of a quantized base model. Initial experiments indicate that quantization down to 8 bits performs well, even for models originally trained with \texttt{bfloat16} precision.

\end{itemize}

\subsection{Datasets}

Each adapter is trained using specialized datasets, summarized in Table~\ref{tab:table_loradata} along with citations to their sources and original references.

\subsection{Training strategy} Training is performed in multiple stages. We start with the Zephyr-7B-$\beta$ model~\cite{HuggingFaceH4/zephyr-7b-betaFace}, which was developed based on the Mistral-7B model~\cite{Jiang2023Mistral7B}, and proceed to train a series of adapters using the datasets outlined previously. The Zephyr-7B-$\beta$ tokenizer is used throughout.

The initial eight adapters are trained using question-answer pairs provided by the datasets (see Table~\ref{tab:table_loradata}).

The protein mechanics tasks follow a similar training approach, but utilize specialized forward and inverse instruction sets. The bidirectional instruction tasks include:

\begin{quote}
\tiny {\texttt{CalculateForce< G E E C D C G S P ....> [0.310]\\
CalculateEnergy< G E E C D C G S P ....> [0.220]\\ 
CalculateForceHistory< G E E C D C G S P ....> [0.004,0.034,0.125,0.142,0.159,0.102,0.079,0.073,0.131, ...]\\
GenerateForce<0.310>\,[ G E E C D C G S P ....]\\
GenerateEnergy<0.220>\,[ G E E C D C G S P ....]\\
GenerateForceHistory<0.004,0.034,0.125,0.142,0.159,0.102,0.079,0.073,0.131, ...>\,[ G E E C D C G S P ....] }}
\end{quote}

These correspond to three distinct tasks: determining the maximum unfolding force of a protein, calculating the unfolding energy, and obtaining the force-deformation history. Each of these forward tasks has a corresponding inverse task that produces a protein sequence as output, resulting in a total of six tasks.

Each LoRA adapter has a rank of $r=16$ with a scaling factor $\alpha=16$. The adapters are trained on five specific target modules: \texttt{v\_proj}, \texttt{k\_proj}, \texttt{q\_proj}, \texttt{gate\_proj}, and \texttt{o\_proj}, with a dropout rate of 0.05. Typically, LoRA adapters are trained with low ranks ranging from 4 to 32 because prior research has demonstrated that increasing the rank does not generally enhance performance, but it does raise computational costs. Lower ranks lead to more compact models with fewer parameters~\cite{hu2021lora}. The selections of $r$, $\alpha$, and other parameters follow standard values frequently adopted in the literature that have proven effective. Furthermore, previous experiments conducted by the author of this work tested various configurations and concluded that the chosen values here represent a reasonable compromise~\cite{Luu2023BioinspiredLLM:Materials}.

The X-LoRA scaling head is trained using approximately 2,000 samples, which are a subset of the original dataset employed for training the adapters (we also enhanced the training set with additional multiple-choice questions generated by GPT-4). We utilize a \texttt{paged\_adamw\_8bit} optimizer with a gradient clipping threshold of 0.3, a learning rate of $2 \times 10^{-4}$ accompanied by warmup, and four gradient accumulation steps, using a batch size of one. The training of the X-LoRA model spans roughly 10,000 steps, with the final loss settling around 0.28.

Given that the fundamental Mistral architecture serving as the base model has 32 layers, and adapters are applied in 5 layers of each transformer layer, there are a total of 160 LoRA layers, each replaced by an X-LoRA layer. The X-LoRA scaling head predicts scaling coefficients for each of these 160 LoRA layers. With nine LoRA experts, this leads to a total of $160 \times 9 = 1440$ $\lambda_i$ values computed per token. The model employs a single linear layer followed by ReLU activations and a dropout rate of $p=0.2$, containing approximately 5 million parameters. The implementation supports flexible architectures, allowing for the design and testing of more sophisticated deep classifier heads in other use cases.

Both LoRA adapters and the X-LoRA scaling head are trained using next-token prediction along with cross-entropy loss (we estimate the probability that the model assigns to the correct next token in a sequence based on the preceding tokens; the objective during training is to maximize the likelihood of the correct next token as indicated in the training data). Token shifting is applied, where the sequence is shifted by one token to produce the target sequence the model should predict. For instance, if the input sequence is ``Proteins are fundamental to", the target sequence used for training becomes ``are fundamental to biology". Thus, the model attempts to predict the subsequent token in the sequence.

As illustrated in Figure~\ref{fig:hierarchical_concept}, during training of the X-LoRA scaling head, only the scaling head is updated while all other parameters remain fixed. The original Zephyr tokenizer and prompt template are utilized.

\subsection{X-LoRA-Gemma model} The X-LoRA-Gemma model is created similarly to the X-LoRA model described above, but it is based on the Gemma-7B-it model~\cite{Google/gemma-7b-itFace}, employing four adapters (details are provided in the main text). These four adapters are trained on datasets covering mechanics and materials, protein mechanics, bioinspired materials, as well as an additional quantum mechanics-based molecular properties dataset QM9 (refer to Table~\ref{tab:table_QM9})~\cite{Ruddigkeit2012EnumerationGDB-17,Ramakrishnan2015ElectronicSpace}.

Each LoRA adapter has a rank of $r=16$ and a scaling factor of $\alpha=16$. They are trained targeting seven specific modules: \texttt{q\_proj}, \texttt{o\_proj}, \texttt{k\_proj}, \texttt{v\_proj}, \texttt{gate\_proj}, \texttt{up\_proj}, and \texttt{down\_proj}, with a dropout rate of 0.05. The rank for each adapter is set at $r=16$ with $\alpha=16$. Given four adapters and seven targeted modules, and considering that the Gemma-7B base model has 28 layers, a total of $28\times 28=784$ $\lambda_i$ values are computed per token.

The initial two adapters (bioinspired and mechanics of materials) used to build this model are trained on question-answer pairs from the corresponding datasets (see Table~\ref{tab:table_loradata}).

Training for the protein mechanics tasks follows the approach described earlier, involving bidirectional forward and inverse instruction sets. Additionally, an adapter is trained using the QM9 dataset, comprising two tasks:

\begin{quote}
\tiny {\texttt{CalculateMolecularProperties<CC(C)N1CC1C=O> [0.098,0.358,0.581,0.395,0.309,0.330,0.570,0.649,0.519,0.519,0.519,0.518]\\
GenerateMolecularProperties<0.098,0.358,0.581,0.395,0.309,0.330,0.570,0.649,0.519,0.519,0.519,0.518> [CC(C)N1CC1C=O]\\ 
...
}}
\end{quote}

These correspond to a total of two tasks (predicting molecular properties, the forward task, and the inverse task of generating a molecule with desired properties). All 12 properties included in the QM9 dataset are either predicted or designed for, respectively (refer to Table~\ref{tab:table_QM9}).  
This model contains a total of four adapters.

The X-LoRA scaling head for this model consists of three layers, with the intermediate layer sized at $3072 \times {FF}_{\rm fac} = 12,288$, where the feed-forward multiplier ${FF}_{\rm fac} = 4$. We employ linear layers with ReLU activation functions and a dropout rate of $p=0.2$, totaling approximately 47 million parameters.

The X-LoRA-Gemma scaling head is trained on roughly 17,000 samples. These samples represent a subset drawn from the training dataset used for adapter training; it is important to note that although the X-LoRA-Gemma model only includes four adapters, we utilized samples spanning the entire training set of the aforementioned model plus additional samples from the QM9 dataset. This approach was taken to explore whether the X-LoRA model could acquire further capabilities during the scaling head training phase. Given that the loss converged well to 0.58 by the end of training, along with strong performance on tasks the model was not explicitly trained for (e.g., chemistry tasks demonstrated in the molecule design example presented in the paper), we conclude that this method can generally produce effective outcomes.

We employed a \texttt{paged\_adamw\_8bit} optimizer with gradient clipping set to 0.3, a learning rate of $1\times 10^{-4}$ accompanied by a warmup period, and four steps of gradient accumulation, using a batch size of one. The X-LoRA-Gemma model was trained for approximately 62,000 steps.

The original Gemma tokenizer and prompt template were utilized.

\subsection{Adversarial agentic modeling} Our approach to adversarial agentic modeling involves creating two X-LoRA agents. One agent specializes in generating questions: it initiates the first query and then progressively formulates further questions aimed at uncovering challenges and weaknesses in the answers. The other agent acts as the respondent, continuously providing replies. This methodology is implemented using \{Guidance\}~\cite{Guidance-ai/guidance:Models.}.

For the example connecting music and mechanics mentioned in the text, the roles of the two agents are outlined as follows. First, the physicist who poses questions:

\begin{quote}
\small\texttt{You are a critical physicist. You participate in a discussion from a physics standpoint. Keep your responses concise, and consistently challenge statements in a provocative manner.

As a creative thinker, you bring in ideas from other disciplines and push conceptual boundaries.}
\end{quote}

Second, a philosopher (in some cases this role is adapted to represent a biology expert):

\begin{quote}
\small\texttt{You are a philosopher knowledgeable in biology, chemistry, and mathematics. You engage in the discussion.

Provide brief but precise and inventive answers. You generate excellent ideas and novel lines of thought, always maintaining logical consistency.}
\end{quote}

The question-asking agent is explicitly directed to reflect on the question and preceding answers, and to formulate new incisive questions. This is achieved through a prompt structured as follows:

\begin{quote}
\small\texttt{Examine the following question and answer.

\#\#\# Question: <question>

\#\#\# Response: <response>

\#\#\# Instruction: Provide a SINGLE follow-up question that critically examines the response.

Do NOT answer or comment on the question yet.

The single question is: }
\end{quote}

After the dialogue has continued for $N$ turns, the model's task is to 1) summarize the main insights discussed, 2) enumerate the key insights as bullet points, and 3) determine the most significant conclusion.

The raw transcript of the conversation serves as the foundation for subsequent analysis through graph construction.

\subsection{Knowledge graph generation} We employ Zephyr-7B-$\beta$ to extract triplets from the text, using the method described in~\cite{Rahulnyk/knowledge_graph:QnA} with enhancements inspired by the Llama Index graph generation algorithm~\cite{Liu_LlamaIndex_2022}. The graphs are visualized using NetworX\cite{Networkx/networkx:Python} and Pyvis~\cite{WestHealth/pyvis:Graphs.}. A principal directive provided, following the approach recommended in~\cite{Liu_LlamaIndex_2022}, is:

\begin{quote}
\small\texttt{Your task is to extract the ontology of terms mentioned in the given context. These terms should capture the core concepts according to the context.

Present your output as a JSON list. Each list element should include a pair of terms and the relation between them, for example: [...] }
\end{quote}

We instruct the construction of triplets consisting of nodes and edges as follows, consistent with the method in\cite{Liu_LlamaIndex_2022}:

\begin{quote}
\small\texttt{\begin{itemize}
            \item "node\_1": "A concept extracted from the ontology"
            \item "node\_2": "A concept related to node\_1 within the ontology"
            \item "edge": "A concise description of the relationship linking node\_1 and node\_2"
        \end{itemize}
        }
\end{quote}

We also provide examples of ontological analysis to the model, such as:

\begin{quote}
\small\texttt{
       Context: ```Spider silk is a strong natural fiber used to catch prey in a web.``` }
\end{quote}

Then, example triplets are given, for instance:

\begin{quote}
\small\texttt{\begin{itemize}
            \item "node\_1": "spider silk"
            \item "node\_2": "fiber"
            \item "edge": "is"
        \end{itemize}
        }
\end{quote}

The collected graph data is subsequently organized and divided into communities via the Girvan-Newman algorithm~\cite{Girvan2002CommunityNetworks}.

\subsection{Visualization of molecular structures}

We utilize PyMOL~\cite{Schrodinger2015The1.8} to depict and analyze the predicted protein structures. Various scripts and functions within the software help identify features of the proteins, such as secondary structure elements, hydrophobic/hydrophilic regions, disulfide bridges, and hydrogen bonds.

\section*{Data Availability Statement} The code and data supporting this study's results are publicly accessible at \url{https://github.com/EricLBuehler/xlora} and \url{https://huggingface.co/lamm-mit/x-lora}. Additional datasets employed in this work are cited (see specifically Table~\ref{tab:table_loradata}).

\end{document}